\chapter{CRONOGRAMA Y PRESUPUESTO DE IMPLEMENTACIÓN DE LA MÁQUINA}
\label{cap:cronograma-presupuesto}

El alcance de esta tesis es el diseño de la máquina por cálculo, simulación y modelado CAD; su fabricación, montaje y ensayo no forman parte del proyecto. Por ello, este capítulo presenta el presupuesto de inversión de la máquina diseñada en el Capítulo~\ref{cap:diseno}, como estimación de clase 4, y lo compara con el costo de la limpieza manual descrita en el Capítulo~\ref{cap:introduccion}. El cronograma del desarrollo académico se presenta en la Sección~\ref{sec:cronograma-proyecto} del Anexo~\ref{anx:portafolio}.

\section{Identificación de actividades y su duración para la implementación}

El presupuesto se organiza en las cuatro actividades de la Tabla~\ref{tab:actividades-costos}. Las tres primeras se ejecutaron en el proyecto; la cuarta corresponde a la etapa de implementación.

\begin{table}[htbp]
\centering\small\singlespacing
\caption{Actividades del presupuesto de implementación}
\label{tab:actividades-costos}
\begin{tabularx}{\textwidth}{@{}p{0.6cm}YP{3.8cm}P{2.6cm}@{}}
\toprule
\textbf{ID} & \textbf{Actividad} & \textbf{Resultado} & \textbf{Estado}\\
\midrule
1 & Metrado de materiales y horas de taller a partir del modelo CAD integrado & Masas por pieza y horas por complejidad & Ejecutada\\
2 & Consulta de precios de mercado y conversión a soles & Base de precios del Anexo~\ref{anx:calculos} & Ejecutada\\
3 & Selección comercial de bomba, boquillas, accionamiento, filtro, tanques y tablero & Componentes de la Tabla~\ref{tab:seleccion-componentes} & Ejecutada con catálogos\\
4 & Cotizaciones formales de las diez partidas de mayor costo & Presupuesto con precios cotizados & Etapa de implementación\\
\bottomrule
\end{tabularx}
\end{table}

\section{Diagrama de Gantt}

Las actividades del presupuesto figuran en la fase 3 del diagrama de Gantt del proyecto (Figura~\ref{fig:gantt-2} del Anexo~\ref{anx:portafolio}): la selección final estaba programada del 13 al 17 de noviembre de 2026 y la evaluación económica, del 16 al 19 de noviembre, antes de la entrega del informe final del 22 de noviembre. Ambas forman parte de la ruta crítica de esa fase. La selección se adelantó y se cerró con los catálogos disponibles; las cotizaciones formales quedan fuera del plazo del proyecto.

\section{Identificación de recursos materiales y no materiales}

Los recursos materiales son los componentes de la Tabla~\ref{tab:seleccion-componentes} y las piezas de fabricación propia de los planos del Anexo~\ref{anx:planos}. Se agrupan en ocho partidas: mesa, centrado y traslado; aspersión y accionamiento; hidráulica y filtrado; tanques y calentamiento; cabina; drenaje; tablero eléctrico y control; y obra civil, montaje y puesta en marcha. La Tabla~\ref{tab:recursos-implementacion} resume los recursos con mayor peso en el costo y los no materiales.

\begingroup
\small\singlespacing
\setlength{\tabcolsep}{4pt}
\begin{longtable}{@{}>{\raggedright\arraybackslash}p{3.4cm}>{\raggedright\arraybackslash}p{\dimexpr\textwidth-3.4cm-3.4cm-6\tabcolsep\relax}>{\raggedright\arraybackslash}p{3.4cm}@{}}
\caption{Recursos materiales y no materiales}\label{tab:recursos-implementacion}\\
\toprule
\textbf{Recurso} & \textbf{Designación y cantidad} & \textbf{Origen}\\
\midrule\endfirsthead
\caption[]{Recursos materiales y no materiales (continuación)}\\
\toprule
\textbf{Recurso} & \textbf{Designación y cantidad} & \textbf{Origen}\\
\midrule\endhead
\midrule\multicolumn{3}{r}{\textit{Continúa en la página siguiente.}}\\\endfoot
\bottomrule\endlastfoot
\multicolumn{3}{@{}l}{\textit{Materiales y equipos}}\\
Bastidor de la mesa & W6$\times$9, ASTM A36, con uniones soldadas; 168 kg & Fabricación en taller\\
Pedestales del riel & 24 pedestales de acero inoxidable 304; 701 kg; 96 anclajes M16 & Fabricación en taller\\
Rieles y ruedas & Dos rieles A45 de 6540 mm; cuatro ruedas Blickle SPKVS 200K & Comercial \cite{arcelor2023,blickle2026}\\
Traslado & Dos cadenas 08B de inoxidable, reductor sinfín y volante & Comercial y taller\\
Cabina & Paneles de acero inoxidable 304 de 1,5 mm; 768 kg & Fabricación en taller\\
Soporte superior & Postes, vigas, bases y cartucho del eje; 559 kg & Fabricación en taller\\
Bomba y filtro & Grundfos CRN 15-3 de 316 con variador; filtro de bolsa Pentair ES2 & Comercial\\
Boquillas y tubería & 24 boquillas VeeJet de 316; tubería de 316L Sch10S & Comercial \cite{manSpray,asmeb3619}\\
Accionamiento del distribuidor & Motorreductor SEW SF47R17DRN63MS4 y limitador Mayr ROBA-slip & Comercial\\
Tanques y calentamiento & Tanques de lavado (600 L) y de enjuague de inoxidable 304, con calentamiento y dosificación & Estimación preliminar\\
Agente de limpieza & Zep Recirculating Detergent, alcalino de baja espuma, 7--15 g/L & Comercial \cite{zep36342}\\
\multicolumn{3}{@{}l}{\textit{No materiales}}\\
Mano de obra & Taller de calderería y mecanizado a S/~47,5/h, con sueldos de Arequipa; dos técnicos durante 15 días para montaje y puesta en marcha & Tarifa estimada\\
Programas & SAP2000, ANSYS Fluent 2026 R1 (licencia Student), FreeCAD, SolidWorks, CalculiX, OpenFOAM y Python & Disponibles\\
\end{longtable}
\endgroup

El detergente es un producto en polvo para lavadoras de piezas por aspersión con recirculación; no forma espuma por encima de 61~$^\circ$C, admite hasta 70~$^\circ$C y es compatible con acero, cobre y latón \cite{zep36342}, de modo que cubre los aros, los rodillos y la jaula a 62~$^\circ$C. A 7--15 g/L, la solución conserva las propiedades del agua usadas en el cálculo hidráulico y en el modelo de mojado.

\section{Determinación de costos unitarios y totales}
\label{sec:costos-unitarios}

Las piezas fabricadas se costearon con la masa de cada pieza en el modelo CAD por el precio del material, más las horas de taller por una tarifa de S/~47,5/h, obtenida del sueldo de un soldador en Arequipa con cargas sociales, gastos y utilidad; las horas por kilogramo dependen de la complejidad de la pieza y son una hipótesis. Los componentes comerciales llevan su precio publicado o un precio estimado a partir de productos similares, con un rango más amplio; lo importado lleva un factor de 1,20 por flete, seguro y aranceles. Se usaron S/~3,448 por dólar y S/~3,91 por euro, y la fuente de cada precio figura en la Sección~\ref{sec:base-precios} del Anexo~\ref{anx:calculos}. A cada partida se le asignó un rango de 0,85 a 1,20 si tiene fuente y de 0,70 a 1,50 si es una estimación. La Tabla~\ref{tab:costos-componentes} muestra las diez partidas de mayor costo, que suman la mitad del subtotal.

\begin{table}[htbp]
\centering\small\singlespacing
\caption{Costos unitarios y totales de las diez partidas de mayor costo}
\label{tab:costos-componentes}
\begin{tabularx}{\textwidth}{@{}Yrrl@{}}
\toprule
\textbf{Partida} & \textbf{Cantidad} & \textbf{Total (S/)} & \textbf{Base del precio}\\
\midrule
Paneles de la cabina, inoxidable 304 de 1,5 mm & 768 kg & 23\,999 & Material con fuente; horas estimadas\\
Pedestales del riel, inoxidable 304 & 24 & 21\,895 & Material con fuente; horas estimadas\\
Bomba Grundfos CRN 15-3 de 316, motor de 4 kW & 1 & 17\,388 & Estimado de la misma serie\\
Montaje, pruebas y puesta en marcha & 240 h & 11\,399 & Estimado\\
Dosificadoras de detergente e inhibidor & 2 & 10\,127 & Precio publicado\\
Soporte superior con bases y cartucho & 559 kg & 10\,067 & Material con fuente; horas estimadas\\
Filtro de bolsa Pentair ES2 de 316L & 1 & 9138 & Precio publicado\\
Obra civil: pozo, piso húmedo, colector y bordillos & 1 & 8500 & Estimado\\
Ruedas Blickle SPKVS 200K & 4 & 7448 & Estimado\\
Traslado por doble cadena con volante & 1 & 6885 & Mixto\\
\midrule
\textbf{Diez partidas} & & \textbf{126\,846} & 48,7\,\% del subtotal\\
\bottomrule
\end{tabularx}
\par\smallskip\footnotesize
Montos sin IGV, con la tarifa de taller de Arequipa. La bomba CRN 15-3 de 316 con motor de 4 kW \cite{grundfoscr} se valoró con el precio publicado de la Grundfos CR 10-06, de la misma serie, multiplicado por 1,4 por el acero 316 y el tamaño. La obra civil es una estimación preliminar de un elemento que la tesis no diseña.
\end{table}

Cerca del 20\,\% del subtotal tiene un precio publicado; el resto es una estimación o una combinación de material con fuente y horas estimadas. La fabricación suma S/~95\,142, de los cuales S/~39\,534 son mano de obra de taller. La estimación anterior, de S/~33\,855, era unas nueve veces menor porque solo incluía los componentes comerciales.

\section{Formulación del presupuesto de implementación}

La Tabla~\ref{tab:presupuesto-completo} agrupa la estimación por partida. Al subtotal se le suma un 5\,\% de ingeniería de fabricación y gestión, y un 15\,\% de contingencia. El total, sin IGV, es de S/~314\,310, con un rango de S/~244\,124 a 428\,655. No incluye el rodamiento, el galpón, el puente grúa ni la acometida eléctrica. Las partidas de tanques y calentamiento, cabina, drenaje, tablero y obra civil son estimaciones preliminares de elementos que la tesis representa en el CAD o en el esquema hidráulico, pero no dimensiona.

\begin{table}[htbp]
\centering\small\singlespacing
\caption{Presupuesto de inversión de la máquina por partida (estimación de clase 4, sin IGV)}
\label{tab:presupuesto-completo}
\begin{tabular}{@{}lrr@{}}
\toprule
\textbf{Partida} & \textbf{Base (S/)} & \textbf{Rango (S/)}\\
\midrule
Mesa, centrado y traslado & 60\,076 & 42\,666--88\,891\\
Aspersión y accionamiento & 32\,286 & 23\,956--45\,722\\
Hidráulica y filtrado & 40\,100 & 29\,441--57\,407\\
Tanques y calentamiento* & 32\,778 & 25\,190--44\,675\\
Cabina hermética* & 45\,591 & 32\,673--66\,865\\
Drenaje* & 5689 & 3983--8534\\
Tablero eléctrico y control* & 26\,364 & 20\,014--36\,429\\
Obra civil, montaje y puesta en marcha* & 20\,615 & 14\,430--30\,922\\
Ajuste de la mano de obra a la tarifa de Arequipa & $-3201$ & $-4802$ a $-2241$\\
\midrule
Subtotal & 260\,298 & \\
Ingeniería de fabricación y gestión (5\,\%) & 13\,015 & \\
Contingencia (15\,\%) & 40\,997 & \\
\textbf{Total sin IGV} & \textbf{314\,310} & 244\,124--428\,655\\
\bottomrule
\end{tabular}
\par\smallskip\footnotesize
*Estimación preliminar de elementos que la tesis no dimensiona. Las partidas incluyen la mano de obra de taller con una tarifa inicial de S/~50,5/h; la fila de ajuste la lleva a la tarifa de Arequipa, S/~47,5/h, con el mismo criterio de rango.
\end{table}

El beneficio se mide frente a la limpieza manual que la empresa de remanufactura de referencia aplica a cada rodamiento: unos diez turnos-persona, 10 L de solvente y 75 kg de residuos contaminados (Capítulo~\ref{cap:introduccion}). Con el sueldo de un técnico en mantenimiento en Arequipa, de S/~1289 a 2927 mensuales y S/~1942 en promedio (Sección~\ref{sec:base-precios} del Anexo~\ref{anx:calculos}), y 26 turnos al mes, la mano de obra cuesta S/~496--1126 por rodamiento, S/~747 con el promedio y sin cargas sociales. Con el valor actual neto, recuperar S/~314\,310 en cinco años al 10\,\% solo con ese ahorro exigiría reemplazar el 925\,\% de la limpieza manual con doce rodamientos al año. Ese valor actual neto no cierra con ningún volumen razonable, de modo que la evaluación se plantea como el beneficio por rodamiento que haría falta para recuperar la inversión:
\begin{equation}
  B^{*}=\frac{I}{A\,n},\qquad A=\frac{1-(1+i)^{-H}}{i},
  \label{eq:beneficio-requerido}
\end{equation}
donde $I$ es la inversión, $n$ el número de rodamientos por año y $A$ el factor de valor actual de una anualidad. Con un horizonte de $H=5$ años y una tasa del 10\,\% anual ($A=3{,}791$), ambas hipótesis de evaluación, la Tabla~\ref{tab:beneficio-requerido} da el resultado.

\begin{table}[htbp]
\centering\small\singlespacing
\caption{Beneficio por rodamiento necesario para recuperar la inversión en cinco años al 10\,\% anual}
\label{tab:beneficio-requerido}
\begin{tabular}{@{}rrrr@{}}
\toprule
\textbf{Rodamientos por año} & \textbf{Base (S/)} & \textbf{Rango (S/)} & \textbf{Mano de obra manual (S/)}\\
\midrule
4 & 20\,729 & 16\,100--28\,270 & 747\\
12 & 6910 & 5367--9423 & 747\\
24 & 3455 & 2683--4712 & 747\\
48 & 1727 & 1342--2356 & 747\\
\bottomrule
\end{tabular}
\end{table}

La mano de obra cubre entre el 4 y el 43\,\% de ese beneficio, según el volumen. El resto tendría que venir del solvente que se deja de usar (10 L por rodamiento), de los residuos peligrosos que no se generan (75 kg), de la menor exposición del personal y, sobre todo, del tiempo de parada del molino que se acorte. Esos datos dependen del taller y de cada planta, y no forman parte de esta tesis. En el Perú operan unos 120 rodamientos 241/900 \cite{taller_ingeniero_2026}, por lo que la viabilidad económica depende de que el taller concentre la remanufactura de varias plantas.

El costo de operación por ciclo comprende la energía y el detergente. Llevar el rodamiento de 15 a 60~$^\circ$C requiere del orden de 17 kWh ($Q=mc\,\Delta T$, con $c=0{,}46$ kJ/(kg\,K) para el acero), que dominan el consumo; a S/~0,7238/kWh (tarifa BT5B de referencia, con IGV), cuestan unos S/~12,6 por rodamiento. Como referencia, una lavadora comercial de tamaño similar instala 60 kW de calentamiento eléctrico \cite{stingray_6075}. El detergente cuesta S/~31,99/kg, es decir, S/~237--473 por metro cúbico de baño a 7,4--14,8 kg/m$^3$. El agua de reposición depende del arrastre y de las purgas, que no se modelaron. La restricción presupuestal RS-10 no tiene un techo acreditado; el total queda como referencia para las cotizaciones formales.
